\section{Finite numerical queries with bounded suspended storage}
\label{sec:arith}

The preceding sections define finite estimator tables and bound their
statistical errors.  We now compute digits of those estimator values.
For a fixed table tuple, including its statistical budget, the structural
conventions are fixed and each entry has one exact mathematical value.
Increasing the digit position requests a later digit of that same value.
Each request is a finite deterministic computation.

A query returns one bounded digit.  During its evaluation, the space
needed by a suspended caller must be bounded by the size of its
structural descriptors and a constant times the decrease from its
output digit position to the requested operand digit position.  An operand digit position may
exceed the output digit position by an amount proportional to descriptor size;
the structural term also pays for this overshoot.  These differences
will telescope when arithmetic queries are nested inside graph procedures.
Thus the relevant question is which data remain live at each operand request.

\subsection{Streams, records, and uniformity}

We use integer digit positions $\nu,j\geq 0$.  Known binary scales are separated
from numerical values, so a normalized value belongs to $[-1,1]$.
The fixed digit alphabet is
\[
 \mathcal A_{\rm dig}=\{-6,-5,\ldots,5,6\}.
\]
Every digit is encoded by a fixed number of bits.  An operand stream
$(x_j)_{j\geq0}$ denotes the exact real number
\begin{equation}
 x=\sum_{j\geq0}x_j2^{-j},\qquad x_j\in\mathcal A_{\rm dig}.
 \label{eq:arith-stream}
\end{equation}
In particular, for its truncation $x^{[K]}=\sum_{j=0}^Kx_j2^{-j}$,
\begin{equation}
 |x-x^{[K]}|\leq 6\,2^{-K}.
 \label{eq:arith-tail}
\end{equation}
Redundant signed-digit representations have a long history in
arithmetic~\cite{Avizienis1963}.  The redundancy lets consecutive
finite approximations define an exact stream even when they use
different truncations.

An arithmetic root has public parameters, its digit position, and a description
of its operand interfaces.  We require these parameters to be
regenerable from the outer controller records without operand or graph
queries.  They are fixed inputs to each numerical computation; their
full binary encodings are not copied at every gate.  An operand label
specifies its interface and descriptor; it never specifies an unknown
vertex address from a previous cursor position.

The following lemma reduces each digit query to a fixed number of
residue bits.

\begin{lemma}[Exact streams from approximation residues]
\label{lem:arith-residue-stream}
Suppose $|z|\leq1$ and, for each $\nu\geq0$, a deterministic integer
$N_\nu(z)$ satisfies
\begin{equation}
 |N_\nu(z)-2^\nu z|<2.
 \label{eq:arith-approximation}
\end{equation}
Set
\begin{equation}
 z_0=N_0(z),\qquad z_\nu=N_\nu(z)-2N_{\nu-1}(z)\quad(\nu\geq1).
 \label{eq:arith-difference-digit}
\end{equation}
Then $z_\nu\in\mathcal A_{\rm dig}$ and these digits represent $z$ exactly.
For any fixed integer $q\geq4$, each output digit is determined by
the required $N_\nu$ residues modulo $2^q$.
\end{lemma}

\begin{proof}
Writing $e_\nu=N_\nu-2^\nu z$ gives
$z_\nu=e_\nu-2e_{\nu-1}$ for $\nu>0$, hence $|z_\nu|<6$; the weaker alphabet
bound $|z_\nu|\leq6$ is convenient throughout.  At place zero,
$|N_0|<3$, so $|N_0|\leq2$.  The partial sum through place $T$ is
$2^{-T}N_T$, which tends to $z$ by Equation~\eqref{eq:arith-approximation}.
Finally, reduction modulo $2^q$ is injective on the possible signed
digits because $2^q>12$.  Thus the signed difference can be recovered
from the two residues without forming the full integers $N_\nu$.
Concretely, for $\nu>0$ compute
$(N_\nu\bmod 2^q)-2(N_{\nu-1}\bmod 2^q)$ modulo $2^q$, then use the
fixed finite table which sends each occurring residue to its unique
representative in $\mathcal A_{\rm dig}$.  At $\nu=0$, decode the residue of
$N_0$ in the same way.  This final decoding has constant storage and
introduces no comparison of unknown exact real numbers.
\end{proof}

The integer $N_\nu$ is always defined before taking its residue.
Different places may use different truncations or finite circuits.
Only determinism and Equation~\eqref{eq:arith-approximation} are
required; consistency between the finite approximations is supplied
by Equation~\eqref{eq:arith-difference-digit}.  In particular, when a
digit at $\nu$ uses $N_{\nu-1}$, the latter is evaluated by its own
place-$(\nu-1)$ definition.

We compute the required residues by finite Boolean circuits.  To bound
what survives an operand query, we evaluate those circuits from paths
of child choices and regenerate their absolute addresses as needed.
Here $d\geq1$ measures the size of the immediate structural descriptors;
in the application it is the current group's slack.

The local record of a suspended numerical query contains every local bit
that must be retained across an operand query: its binary gate path,
return modes, pending bounded results, stored descriptor choices, and
any explicitly specified retained word.  The input and master choice,
the common vertex cursor, and one reusable scratch area are shared
facilities.  Temporary gate addresses, digit places and scan counters
may occupy scratch while they are being computed.  Any such data needed
after a recursive graph query must either be in the charged local
record or be reconstructible from the root parameters and that record.
An operand-dependent word retained across the query is charged even if
it was originally formed in scratch.

\subsection{The numerical interface}

We state the composition result with its syntactic hypotheses.
They concern finite arithmetic recipes.

Application to graph procedures requires a separate certificate.  The
semantic descriptor domains and structural validity conventions are
fixed independently of the requested digit position; selected operands
already have defined digit interfaces; reciprocal arguments lie in the
stipulated positive intervals; and root parameters remain recoverable
after lower calls.  Digit-position-dependent truncations and comparison circuits
may change $N_\nu$, but not its target recipe value.
Lemmas~\ref{lem:ctrl-recipes} and~\ref{lem:ctrl-replay} verify these
conditions for the estimator procedures.  An incoming order statistic over the
whole vertex universe has no prescribed $2^{O(d)}$ descriptor bound and
uses the separate construction of Lemma~\ref{lem:ctrl-score-rank}.

The theorem uses the following terms.  The
\emph{recipe dependency depth} of a recipe is the largest number
of operations on a dependency path through it, ending at an already
supplied operand interface and counting each indexed sum, average or order statistic as a single
operation whatever the depth of its internal tree.  The internal depth
of an indexed sum, average or order statistic is accounted for separately by the order-statistic and
sum lemmas below.

A \emph{discrete operand} is a supplied operand of a recipe that returns
a discrete domain, validity or branch decision, a bit or a bounded
symbol.  Its definition does not depend on the requested numerical
digit position, and it is requested at digit position zero.  The theorem bounds the
records retained above this request; the implementation of the operand
is a separate obligation.
A \emph{comparison with a fixed margin} is a discrete decision computed inside a recipe
by a finite-margin comparison with known operand magnitude bounds
$2^{\kappa d}$ and prescribed buffer widths at least $2^{-\kappa d}$,
for a fixed $\kappa$.  It uses $O(d)$ normalized digit places, with
approximation and tie conventions fixed from the structural parameters,
independently of the requested digit position; its structural starting
digit position is zero.  The test of Lemma~\ref{lem:arith-buffered-test}, with a
threshold $u\ge\rho\ge2^{-\kappa d}$, accepts every value greater than
$2u$ and accepts only values greater than $u$.

A \emph{public coefficient leaf} is a leaf whose normalized numerical
value is determined by the public root parameters alone, without any
operand value; its digit procedure must satisfy hypothesis~8 below.
Known rate and field coefficients and walk probabilities are examples
(Remark~\ref{rem:arith-clocks}), whereas the coefficients formed in
Lemma~\ref{lem:arith-reciprocal} from an operand prefix are not.

The \emph{suspended local record} at an operand request is the local
record retained across that request, as described above.  The
\emph{intrinsic local record length} is the total length, at any instant,
of records used to evaluate the recipe itself, including its internal
arithmetic calls and retained words.  Its maximum therefore counts both
active and suspended records within the recipe.  It excludes the shared
facilities described above and records inside external operand,
structural or graph procedures, whose storage requires the separate
graph certificate.

\begin{theorem}[Uniform numerical interface]
\label{thm:arith-interface}
Fix the digit alphabet, positive reciprocal intervals, and
constants bounding the following recipe class.  A recipe at
slack $d\geq1$ satisfies the following conditions.
\begin{enumerate}
\item It has a fixed number of stages built from bounded
arithmetic operations, fixed piecewise formulas using maxima,
minima and clipping, reciprocals on the fixed positive compact
intervals, indexed finite order statistics, and sums or averages followed by known
rescalings and bounded clipping.
\item Each order statistic has a public semantic list
length and an $O(d)$-bit rank index; absent descriptors have prescribed
constant dummy values, and out-of-range requests return a prescribed
normalized constant.
\item An empty sum is zero; an average has a public
nonempty domain.
\item Binary scale exponents have magnitude $O(d)$, and
terms and intermediate outputs have known range bounds $2^{O(d)}$.
\item The stage and descriptor
descriptions are generated from the root parameters by index
computations using $O(B+d+\nu)$ scratch and no operand queries, and the
total recipe dependency depth is bounded independently
of $d$.
\item Each indexed domain has size $2^{O(d)}$ and descriptors
of length $O(d)$.
\item Discrete domain and validity decisions are
discrete operands or comparisons with fixed margins with their own fixed
precision conventions.
\item Public coefficient leaves have exact
normalized digit procedures using $O(B+d+\nu)$ shared scratch
and making no operand or recursive graph queries.
\end{enumerate}

Then each normalized recipe value has a deterministic exact
$\mathcal A_{\rm dig}$-stream.  For its digit at digit position $\nu$ there are uniform
constants $a,b,c$ such that every requested operand digit $j$
satisfies
\begin{equation}
 j\leq \nu+cd,\qquad
 \text{total suspended local binary record length}
       \leq ad+b(\nu-j).
 \label{eq:arith-export}
\end{equation}
An output code bit satisfies the same bounds.  A discrete operand
is requested at digit position zero and has suspended local length at
most $ad+b\nu$.  The maximum intrinsic local record length is
$O(d+\nu)$.

All gate, digit and descriptor addresses can be regenerated from
root parameters and the saved local records using
$O(B+d+\nu)$ shared scratch and no recursive graph calls.  When
$d,\nu=O(B)$ this is $O(B)$ scratch.  If a public rational
coefficient has $O(B)$-bit numerator and denominator that can
be generated in $O(B)$ scratch, its queried digits in this digit position
range also require only that shared $O(B)$ scratch.
\end{theorem}

The proof in Section~\ref{subsec:arith-interface-proof} combines the
primitive record bounds below. Their common coefficient on the digit position
decrement makes those decrements telescope; the multiplication and
reciprocal constructions establish this bound for their long intermediate
calculations.

\subsection{Evaluation from regenerated addresses}

In particular, reconstructing an absolute digit place may use a binary
counter in scratch, while the suspended record stores only the choices
needed to recover that place.  The next lemma makes this distinction
precise for the finite circuits used below.

\begin{lemma}[Evaluation from regenerated gate addresses]
\label{lem:arith-gate-evaluation}
Consider a finite Boolean circuit of bounded fan-in, with a specified
output gate.  Conditional gates may evaluate their selector first and
then only their selected branch.  Suppose the type and immediate predecessors of a gate,
or its operand label and digit place if it is an input gate, can be
computed from the root parameters and a path of child choices from
the output gate.  Suppose these computations use $O(B+\nu+d)$ scratch
bits and make no operand or graph queries.  Then depth-first
evaluation retains $O(1)$ bits per gate on its current path, and
requires no saved absolute gate address per gate.  Its suspended
record above an operand input has length at most a constant times
the length of that path, plus any explicitly retained root-local
record.
\end{lemma}

\begin{proof}
At a gate, save its next-child or return mode and the already computed
Boolean input, if any.  At a conditional gate also save the selected
branch bit.  The bounded fan-in makes this a constant-size
record.  To locate the next child, scan the saved path from its root
and regenerate successive gate descriptions in scratch.  The final
description determines the next child choice.  On return, its Boolean
result is combined with the saved bit and the local record is popped.
All temporary addresses and scan counters are discarded before an
operand query; a later scan recreates them from the same root and
saved choices.  Circuit acyclicity gives termination.  This is an
evaluation of the circuit unfolded into a formula; repeated gates are
recomputed, so no array of previously computed values is stored.
\end{proof}

\subsection{Truncations, comparisons, and binary arithmetic}

\begin{lemma}[Finite comparisons and bounded elementary operations]
\label{lem:arith-elementary}
On normalized input streams, a pair average, maximum, minimum, or
fixed binary shift with normalized output admits deterministic integers
satisfying Equation~\eqref{eq:arith-approximation}.  Their residue
procedures, for any fixed $q$, have the following properties.  There
are constants $C_0,C_1$ such that an operand place $j$ obeys
$j\leq \nu+C_0$, its suspended local record has length at most
$C_0+C_1(\nu-j)$, and the intrinsic local space is $O(\nu+1)$.
Fixed finite compositions of these normalized operations have the
same bounds, with enlarged constants.  Fixed Boolean operations on
returned digit codes, including sign decoding and conditional selection
of code bits, add only constant local overhead.
\end{lemma}

\begin{proof}
First consider comparing two truncations through $K=\nu+h$, where $h$
is a sufficiently large fixed integer.  Their difference digits have
magnitude at most 12.  Scanning from place zero to $K$, update a signed
prefix by $p\leftarrow2p+\delta$.  While $|p|\leq12$, retain $p$
exactly.  Once $p>12$ or $p<-12$, retain only the corresponding
absorbing sign state.  This sign is correct for every finite
continuation: the remaining suffix, in the current scale, has
magnitude strictly less than 12.  The state space is finite, so each
transition has a fixed bounded-fan-in Boolean circuit.  The sign
at the end determines the exact order of the two finite truncations.
The path from that sign bit to an input at place $j$ has length
$O(K-j+1)=O(\nu-j+1)$.

Let $v_\nu$ be the average, maximum, or minimum of the two
truncations, and put $N_\nu=\lfloor2^\nu v_\nu\rfloor$.  Each operation is
1-Lipschitz for the sup norm on its operands.  Consequently
$|v_\nu-v|\leq6\,2^{-K}$, where $v$ is the corresponding exact output.
Choosing $h$ so that $6\,2^{-h}<1/2$ gives
$|N_\nu-2^\nu v|<3/2$.

To obtain $N_\nu\bmod2^q$, only a fixed number of low places of the
selected truncation are needed.  More explicitly, at scale $2^K$,
terms of input place $j\leq \nu-q$ are multiples of $2^{q+h}$.
Modulo this modulus the remaining $q+h$ places suffice.  For an
average with integer input prefixes $A_K,B_K$ at scale $2^K$,
use $N_\nu=\lfloor(A_K+B_K)/2^{h+1}\rfloor$.  Its residue needs
modulus $2^{q+h+1}$ and one additional retained place:
only places $j\leq \nu-q-1$ can be omitted.  The finite number
of remaining digits can be added and shifted by a constant-size
Boolean circuit.  A maximum or minimum first uses the comparison
bit to select these residue bits.  Its long comparison branch has
the path bound proved above, and its residue branch reaches only
places within a fixed distance of $\nu$.  A fixed shift changes these
places by a fixed amount.  Fixed normalization factors are handled
in this way whenever an intermediate range exceeds $[-1,1]$.

The topology of the comparison chain and the low-residue circuits
depends only on the root digit position, the fixed margins, and the gate
path.  Regenerating a place by replaying that path needs no saved
binary place counter at each transition.  Lemma~\ref{lem:arith-gate-evaluation}
therefore gives the suspended-space bounds.  Their longest paths
have length $O(\nu+1)$.

For composition, use one common coefficient of $(\nu-j)$ for the
finite primitive list.  Increasing such a coefficient by $\Delta$
requires increasing its constant term by at most $\Delta C_0$,
because the operand overshoot is at most $C_0$.  With these common
coefficients, the digit position differences telescope along the composed
path.  The finite number of constant terms remains a constant.
\end{proof}

We will also use the following explicit Boolean arithmetic facts.
A carry bit in binary addition is described by a generate/propagate
pair $(g,p)$.  Composition of adjacent intervals is the associative
operation
\[
 (g_2,p_2)\circ(g_1,p_1)
   =(g_2\lor(p_2\land g_1),\ p_2\land p_1).
\]
Balanced trees of these pairs give every carry in $O(\log(W+1))$
depth for $W$-bit words.  A balanced tree adding $R$ such words modulo
$2^W$ has depth
\begin{equation}
 O\bigl(\log(R+1)\log(W+1)\bigr).
 \label{eq:arith-adder-depth}
\end{equation}
The trees can be padded by zero words to powers of two, so their
gate addresses are computable by ordinary index arithmetic.
For counts of $R$ Boolean bits one can do better. In one column,
three bits satisfy
\[
 a+b+c=(a\mathbin\oplus b\mathbin\oplus c)
        +2\bigl((a\land b)\lor(a\land c)\lor(b\land c)\bigr).
\]
Apply this constant-depth circuit independently in every column,
placing the second output in the next column. Thus three equal-width
rows become two rows of the same total integer value; a carry is
not propagated through the columns in this round.  In $O(\log(R+1))$ rounds
only two rows remain.  The word width is $O(\log(R+1))$; a final
carry computation gives total depth $O(\log(R+1))$.  These
constructions use bounded fan-in throughout.

\subsection{Multiplication by minimum-index bands}

Grouping products by their smaller input index leaves a gap between
the output digit position and the requested input digit positions; that gap pays for
the retained circuit path within each band.

\begin{lemma}[Multiplication]
\label{lem:arith-multiplication}
For normalized exact input streams $x$ and $y$, multiplication has all
the conclusions of Lemma~\ref{lem:arith-elementary}.  In particular,
its suspended storage at an operand place $j$ is
$O(1)+c(\nu-j)$, despite the potentially long product expansion.
\end{lemma}

\begin{proof}
Fix $q$ and write $A=6$.  Partition ordered pairs $(i,j)$ of
nonnegative places into a zero band, where $\min(i,j)=0$, and
positive bands
\[
 \mathcal P_D=\{(i,j):D\leq\min(i,j)<2D\},
 \qquad D=1,2,4,\ldots.
\]
Set $h_0=h_\circ$ for the zero band and
$h_D=h_\circ+\lfloor D/4\rfloor$ for positive bands.  The fixed
integer $h_\circ$ will be chosen below.  Retain in band $D$ only
pairs with $i+j\leq \nu+h_D$, and let
\begin{equation}
 S_D=\sum_{\substack{(i,j)\in\mathcal P_D\\i+j\leq \nu+h_D}}
           x_i y_j2^{\nu+h_D-i-j}.
 \label{eq:arith-band-integer}
\end{equation}
For the zero band the same formula uses its defining pair set.
These are mathematical integers.

At any fixed diagonal $i+j=s$, a positive band contains at most
$2D$ ordered pairs.  Indeed the smaller place has $D$ possibilities,
and an off-diagonal pair has two orientations.  The zero band has
at most two pairs.  Thus, in units of $2^{-\nu}$, the total omitted
product terms have absolute sum at most
\begin{equation}
 E_{\mathrm{omit}}
 \leq 2A^2\sum_{D\in\{0,1,2,4,\ldots\}}(D+1)2^{-h_D}.
 \label{eq:arith-band-tail-error}
\end{equation}
The geometric summation of the diagonal weights gives this bound
also when the truncation threshold lies below the first possible
diagonal.  Both the series in Equation~\eqref{eq:arith-band-tail-error}
and $\sum_D2^{-h_D}$ converge.  Choose $h_\circ$ sufficiently large
that their displayed combined error bound, including coefficient
$2A^2$, is less than $1/2$.

There are only finitely many nonempty retained bands.  For a
positive one, $2D\leq \nu+h_D$, whence
\begin{equation}
 D\leq\frac47(\nu+h_\circ).
 \label{eq:arith-last-band}
\end{equation}
The left side of the nonemptiness condition grows monotonically
faster than its right side along the dyadic bands.  We may therefore
stop at the last band satisfying this inequality, retaining empty
bands within the finite prefix if necessary.

Starting with the most distant retained band, define integers
recursively by
\begin{equation}
 F_D=S_D+
       \left\lfloor 2^{h_D-h_{D^+}}F_{D^+}\right\rfloor,
 \qquad
 N_\nu=\left\lfloor2^{-h_0}F_0\right\rfloor,
 \label{eq:arith-band-recursion}
\end{equation}
where $D^+$ is the next band, and the last missing tail is zero.
One internal floor, divided back to scale $2^\nu$, changes the sum
by less than $2^{-h_D}$.  Induction down the bands consequently
shows that the difference between $2^{-h_0}F_0$ and the sum of
all retained terms in scale-$\nu$ units is at most $\sum_D2^{-h_D}$.
Equation~\eqref{eq:arith-band-tail-error} and the choice of
$h_\circ$ make the error before the final floor less than $1/2$.
The final floor costs less than one more unit.  Therefore
Equation~\eqref{eq:arith-approximation} holds for $z=xy$.

We next implement Equation~\eqref{eq:arith-band-recursion} using
residues.  At band $D$ use modulus $2^{q+h_D}$.  An integer term
in Equation~\eqref{eq:arith-band-integer} with $i+j\leq \nu-q$
is a multiple of this modulus and can be omitted from the circuit.
It remains present in the mathematical integer $S_D$.
Only the diagonals
\begin{equation}
 \nu-q<i+j\leq \nu+h_D
 \label{eq:arith-relevant-diagonals}
\end{equation}
need to be processed, giving
$O((D+1)(h_D+q+1))=O((D+1)^2)$ candidate pairs.

Signed floor shifts are well defined on the required residues.
For $h'\geq h$ and $a-b=k2^{q+h'}$,
\begin{equation}
 \left\lfloor\frac{a}{2^{h'-h}}\right\rfloor
 -\left\lfloor\frac{b}{2^{h'-h}}\right\rfloor
       =k2^{q+h}.
 \label{eq:arith-signed-residue-floor}
\end{equation}
This identity holds for all signed integers $a,b$.  Thus a next-band
residue modulo $2^{q+h'}$ supplies the correctly floor-shifted
tail modulo $2^{q+h}$.  Apply the same observation to the final
shift by $h_0$ to obtain $N_\nu\bmod2^q$.
In the finite binary circuit, this shift just discards the lowest
$h'-h$ bits of the canonical $(q+h')$-bit residue.  The displayed
identity proves that this is valid also when the underlying integer
is negative.

For completeness, the finite gate structure can be indexed without
storing an absolute operand place at each gate.  For a positive
band enumerate triples $(r,s,o)$, where
$D\leq r<2D$, $\nu-q<s\leq \nu+h_D$, and $o\in\{0,1\}$.
The triple supplies $(r,s-r)$ if $s-r\geq r$ and $o=0$;
it supplies $(s-r,r)$ if $s-r>r$ and $o=1$; otherwise it supplies
a zero term.  This enumerates the relevant pairs exactly once.
For the zero band use $r=0$ with the same orientation rule.
Negative candidate places are invalid and supply zero.
The offsets $r-D$ and $s-(\nu-q+1)$ have $O(\log(D+2))$ bits.
They and $o$ are obtained from paths in padded binary selection trees.
The full places are then computed from the root digit position in scratch.

Each nonzero term is a constant-bit signed product $x_i y_j$,
shifted into a word of width $W_D=q+h_D=O(D+1)$ and interpreted
modulo $2^{W_D}$.  A bit of this word is a fixed Boolean function
of the two bounded digit codes, the bit index and the known shift;
negative products use their ordinary residue, equivalently sign
extension before reduction.  Add these words and the shifted tail
using Equation~\eqref{eq:arith-adder-depth}.  A band contributes
depth $O(\log^2(D+2))$.  Bit selection for the shifted tail is
wiring and does not require evaluating a full tail word.

The sum of these depths through the bands up to $D$ is
$O(\log^3(D+2))$, hence at most $C(D+1)$ for a fixed $C$.
For an original input digit in a positive band, its partner place
is at least $D$.  Equation~\eqref{eq:arith-relevant-diagonals} gives
\begin{equation}
 j\leq \nu+h_D-D
       \leq \nu+h_\circ-\frac34D.
 \label{eq:arith-multiplication-latency}
\end{equation}
Equivalently, $\nu-j\geq\tfrac34D-h_\circ$, so the full path cost
through these bands satisfies
\[
 C(D+1)\leq C\left(1+\frac43h_\circ\right)
                   +\frac{4C}{3}(\nu-j).
\]
This gives $C_0+C_1(\nu-j)$ with fixed constants.  The zero band
has constant depth and only constant overshoot.  All its gate
descriptions, term coordinates, bit positions and tail shifts are
computable by replaying paths and doing integer index arithmetic
from $\nu$ and $q$.  Lemma~\ref{lem:arith-gate-evaluation} converts these
depth bounds to suspended storage.  Finally,
Equation~\eqref{eq:arith-last-band} bounds every intrinsic path
and local word width by $O(\nu+1)$.
\end{proof}

\subsection{Reciprocals and the early-word convention}

A prefix through roughly half the output digit position makes the remainder
in a first-order reciprocal expansion quadratically small. Releasing
that prefix before later input queries keeps its long word from
surviving a call with little remaining digit position gap.

\begin{lemma}[Reciprocal on a positive compact interval]
\label{lem:arith-reciprocal}
Fix constants $0<a\leq b$.  On inputs $x\in[a,b]$, a reciprocal,
after a fixed known normalization, has the same operand-position,
suspended-record, and intrinsic-space bounds as in
Lemma~\ref{lem:arith-elementary}.  Its local record may contain an
$O(\nu+1)$-bit word only while querying input places
$j\leq\lfloor \nu/2\rfloor+O(1)$.
\end{lemma}

\begin{proof}
The input can first be normalized by a fixed binary scale.  It is
enough to prove the result with a fixed digit alphabet and fixed
positive interval; all constants below may depend on $a,b$.
Choose a fixed power of two $C$ large enough that the desired output
$z=1/(Cx)$ lies in $[0,1]$ and that the coefficients below have
absolute value at most one.  Put
\[
 k=\lfloor \nu/2\rfloor+c,\qquad x^0=x^{[k]},\qquad e=x-x^0,
\]
where $c$ is a sufficiently large fixed integer.  By
Equation~\eqref{eq:arith-tail}, $|e|\leq 6\,2^{-k}$.  Choosing $c$
first ensures $x^0\in[a/2,2b]$ and $|e|\leq1$ at every $\nu$.
Increasing $C$ if necessary gives bounded coefficients
\[
 \alpha=\frac1{Cx^0},\qquad \beta=-\frac1{C(x^0)^2}.
\]
The exact algebraic identity
\begin{equation}
 \frac1{Cx}=\alpha+\beta e+
               \frac{e^2}{C(x^0)^2x}
 \label{eq:arith-reciprocal-taylor}
\end{equation}
shows that the linear expression $T_\nu=\alpha+\beta e$ differs
from $z$ by at most a fixed constant times $2^{-2k}$.
Choose $c$ large enough that
\begin{equation}
 2^\nu|T_\nu-z|\leq\frac14
 \quad\hbox{for every }\nu\geq0.
 \label{eq:arith-reciprocal-remainder}
\end{equation}
Clipping $T_\nu$ to $[-1,1]$ cannot increase this error; denote the
clipped expression by $\widehat T_\nu$.

For this fixed original $\nu$, the coefficients $\alpha,\beta$ are
exact rational numbers.  Their digits are produced by integer
arithmetic, as described below.  The suffix $e$ has the exact stream
with digit zero through $k$ and digit $x_j$ thereafter.  Thus
Lemmas~\ref{lem:arith-elementary} and
\ref{lem:arith-multiplication} provide exact streams for
$\widehat T_\nu$, using only a fixed number of normalized operations.
Explicitly, form the product $\beta e$, then the pair average
$u=(\alpha+\beta e)/2$, clip $u$ to $[-1/2,1/2]$ by a maximum
and a minimum, and finally apply the fixed shift by 2.  Every
exported intermediate value is normalized, and the final value is
$\widehat T_\nu$.  Forming this expression introduces no additional
reciprocal operation.

Let $P_s(\widehat T_\nu)$ be the integer approximation furnished by
those operations at place $s=\nu+3$, and define
\begin{equation}
 N_\nu(z)=\left\lfloor P_{\nu+3}(\widehat T_\nu)/8\right\rfloor.
 \label{eq:arith-reciprocal-approximation}
\end{equation}
The cutoff $k$ and the coefficients remain those of the original $\nu$,
not those of $s$.  Equations~\eqref{eq:arith-approximation} and
\eqref{eq:arith-reciprocal-remainder} give
\[
 |N_\nu(z)-2^\nu z|<1+\frac28+\frac14=\frac32<2.
\]
To obtain this integer modulo $2^q$, ask for the approximation
$P_{\nu+3}$ modulo $2^{q+3}$ and use
Equation~\eqref{eq:arith-signed-residue-floor}.

It remains to justify storage in the coefficient computations.
The signed numerator of $x^0$ in units $2^{-k}$ has $O(k+1)$
bits.  Form it by querying places $0,\ldots,k$ and updating the
integer prefix.  The numerator, a place counter, and all data that
must survive the next query are part of the local record.
To return a coefficient digit of place $r$, numerator squaring and
integer division require $O(k+r+1)$ bits.  Here $r\leq \nu+O(1)$,
because the enclosing expression has fixed depth and constant
overshoot.  After the input prefix is obtained, these integer
computations make no further operand queries.  Standard binary
long multiplication and division, performed by bounded loops with
their partial products or remainder, use only this linear number
of bits.  For $r>0$, computing floors of the coefficient scaled at
$r$ and $r-1$ gives its exact difference digit; at $r=0$ use the
place-zero floor alone, as in Lemma~\ref{lem:arith-residue-stream}.
Clipping is unnecessary
because the coefficients were normalized in advance.

Every original-input query during this formation has
$j\leq k\leq \nu/2+O(1)$.  Hence
\begin{equation}
 \nu+1\leq C_2+C_3(\nu-j)
 \label{eq:arith-early-word-payment}
\end{equation}
for fixed constants, at all such queries.  The enclosing gate path
also has length $O(\nu+1)$, so the complete suspended local record
satisfies the required bound there.  Equivalently, the coefficient
branch is charged against the original reciprocal digit position $\nu$, not
against its possibly small coefficient digit place $r$.

On a suffix branch with $j>k$, the coefficient word is absent.
The Boolean evaluation asks for a coefficient digit at a separate
leaf, returns its bounded code bits, and discards the coefficient
word.  If the other child of a gate subsequently asks for a late
suffix digit, only the bounded pending bits and gate choices are
retained.  The path bound from the preceding arithmetic lemmas
therefore gives $O(1)+c(\nu-j)$ also on this branch.  Repeated
coefficient requests rebuild the word.  The parameter $k$ and the
choice of coefficient are regenerated from the original $\nu$ and
the saved gate path; no earlier cursor address is required.
All intrinsic records have $O(\nu+1)$ bits.
\end{proof}

The coefficients in this proof are operand-dependent, even though they
are rational once $x^0$ is known.  They therefore do not qualify as
public coefficient leaves.  A coefficient request first forms its
charged early prefix, then uses scratch for integer arithmetic, returns
only the requested bounded digit, and releases that prefix before a
suffix branch is evaluated.  The original reciprocal digit position $\nu$ remains
regenerable throughout.  This order of operations is the lifetime rule
used in the suspended-record bound.  In particular, $j\le\lfloor \nu/2\rfloor+c$
gives the explicit payment $\nu+1\le2(\nu-j)+2c+1$ for the early word.

\subsection{Indexed order statistics, sums, scales, and discrete operands}

\begin{lemma}[Indexed finite order statistics]
\label{lem:arith-rank}
Fix $\kappa>0$ and $d\geq1$.  Consider at most $2^{\kappa d}$
normalized values, indexed by descriptors of length at most
$\kappa d$, and an order statistic specified by an $O(d)$-bit
index.  The semantic list length $n$ is public; the requested index
is in $\{1,\ldots,n\}$, or an out-of-range request returns a
prescribed normalized constant.  Absent descriptors may be assigned
prescribed constant dummy values by structural validity predicates.
This order statistic has deterministic integer approximations $N_\nu$
satisfying Equation~\eqref{eq:arith-approximation}.  Their residue
procedures have constants $a,b,c$, depending only on the descriptor
bounds, such that an operand digit query at $j$ obeys
\begin{equation}
 j\leq \nu+cd,\qquad
 L\leq ad+b(\nu-j),
 \label{eq:arith-local-bound}
\end{equation}
where $L$ is the total suspended local record length.  A structural
validity query is assigned digit position zero and has suspended length
$O(d+\nu)$.  The intrinsic local length is $O(d+\nu)$.
\end{lemma}

\begin{proof}
For $1\le k\le n$, the $k$th largest of the semantic $n$-entry list
is its $(n-k+1)$st ascending order statistic, before any extra padding.
Pad to a fixed descriptor domain whose cardinality is a power of
two, adding entries of value $+1$.  These entries do not change
any of the first $n$ ascending order statistics.  An empty list or
an out-of-range request returns its specified default immediately.
Replace invalid original descriptors by the prescribed dummies.
An operand access evaluates validity first and enters the actual
value branch only when valid; an invalid access returns the dummy's
digit.  Thus a dummy does not require a query to an undefined value.
For each valid numerical value use its truncation through
$K=\nu+h$ for fixed large $h$.  Order these finite rational values,
breaking ties by descriptor index.  Compare every pair using the
signed comparison of Lemma~\ref{lem:arith-elementary}.  For each
descriptor count how many descriptors precede it, then compare
that count to the requested rank minus one.  Exactly one
descriptor is selected.  Select its low-residue bits using Boolean
AND and OR trees and finally floor to output scale $2^\nu$.
The selected descriptor is internal to this residue computation.  It
may change with $\nu$; the exact order statistic being approximated is
fixed by the semantic list and the requested index.

If all coordinates of two finite lists differ by at most $\eta$,
their corresponding order statistics differ by at most $\eta$.
Indeed, at least the specified number of the perturbed entries
are no larger than the original order statistic plus $\eta$;
interchanging the lists gives the reverse inequality.  Here
$\eta=6\,2^{-K}$.  Thus the selected truncation, followed by the
final floor, satisfies Equation~\eqref{eq:arith-approximation}
for sufficiently large fixed $h$.  Descriptor ties affect neither
this estimate nor the uniqueness of the finite selection.
This argument perturbs the entire padded exact list.  Some truncated
original entries may exceed $+1$ slightly, but their error, and hence
the error of the selected padded order statistic, is still at most
$\eta$.

Each count has $O(d)$ bits.  Carry-save counting as described
above, the comparison to the requested rank, and selection among
$2^{O(d)}$ descriptors add $O(d)$ Boolean depth above the
pairwise comparison bits and selected residue bits.  A
place-$j$ input to a pairwise comparison consequently has path
length $O(d)+O(\nu-j)$; a residue input is within a fixed distance
of $\nu$ and has path length $O(d)$.  The paths to validity predicates
have length at most the whole circuit depth $O(d+\nu)$, and the
predicates themselves are separate discrete operands of digit position zero.

Each gate is specified by its pair of descriptor indices, its
counting-tree or selection-tree position, and a comparator or bit
position.  The paths determine these indices in $O(d+\nu)$ scratch,
using only the fixed domain and root parameters.  No numerical
comparison is needed to determine the circuit topology.  Apply
Lemma~\ref{lem:arith-gate-evaluation}.  The constants in
Equation~\eqref{eq:arith-local-bound} can absorb the fixed digit
overshoot because $d\geq1$.
\end{proof}

For the compression order statistics, copies of a previously computed default
are ordinary indexed operands: their digit interfaces query the
default's exact stream one bounded digit at a time.  A slot that
chooses between an entry and the default first evaluates its
structural flag and then queries the selected, already-defined stream.
Treating this as a defined operand interface leaves no absent value
for the order-statistic procedure to evaluate.
The required default copies belong to the semantic list; the extra
power-of-two padding in the proof is then supplied by the constant
$+1$.  The structural choice adds only a fixed conditional branch,
so the same digit-position and suspended-record bounds apply.

An order statistic over the full incoming vertex universe need not have a
$2^{O(d)}$ descriptor domain.  That operation uses the terminal
configuration procedure in the next section.

\begin{lemma}[Bounded sums and known scales]
\label{lem:arith-sums-scales}
Suppose an immediate descriptor domain has at most $2^{\kappa d}$
terms, all required binary scale exponents have magnitude at most
$\kappa d$, and every intermediate output has a known range bound
$2^{\kappa d}$.  A sum over an empty domain is zero, and an average
is taken over a public nonempty domain.  Balanced sums or averages of these terms, followed
by the required known rescalings and bounded clipping, satisfy
Equation~\eqref{eq:arith-local-bound} and have intrinsic local
space $O(d+\nu)$.
\end{lemma}

\begin{proof}
Choose one common known power-of-two range bound for the immediate
terms and divide each term by it, so every term lies in $[-1,1]$.
Pad their number to a power of two $2^r$, $r=O(d)$, and use a complete
binary tree of pair averages.  Every internal value stays normalized.
Along one path there are $r$ elementary operations.  Choose common
coefficients in their bounds, as in the last paragraph of the
proof of Lemma~\ref{lem:arith-elementary}.  The constant terms
sum to $O(r)$ and the differences of successive digit digit positions
telescope.  Thus the local length is $O(d)+b(\nu-j)$ and the
total overshoot is $O(d)$.  Restore the common normalization scale
only after the tree of pair averages.  For a sum, combine this factor
with the known padding count $2^r$ in one binary shift of magnitude
$O(d)$.  For an average over an original domain of size $m$, combine
the scale restoration with multiplication of the padded average by
the public factor $2^r/m$.  With minimal padding the latter factor
is in $[1,2)$, so it requires only one normalized public multiplication
and a fixed shift in addition to the common scale restoration.

More explicitly, for a shift by $2^v$, truncate the operand at
$K=\max(0,\nu+v+h)$.  If $\nu+v+h\geq0$, its scaled truncation
error in scale-$\nu$ units is at most $6\,2^{-h}$.  If the scaled
output is known to be normalized, floor in the desired output
scale.  Otherwise, when a bounded clip is required, apply that
clip to the scaled truncation before flooring.  For the
normalized clip to $[-1,1]$, for example, define the approximation
using $\min(1,\max(-1,2^v x^{[K]}))$; its error cannot increase.
This shift-and-clip procedure does not first export an unbounded
intermediate as a normalized stream.
The low residue requires the constant number of places near
$\nu+v$, and a comparison needed for clipping has a constant-state
chain through $K$.  Comparisons to the scaled clipping endpoints
are equivalently comparisons of $x^{[K]}$ to the known endpoints
divided by $2^v$; their finite constants and shifted place data
are public.  When $v\geq0$ these endpoints belong to $[-1,1]$;
when $v<0$ the normalized operand is already in the clip interval
after scaling, so no clipping comparison is needed.  Its path cost is
$O(1)+O(\max(0,\nu+v-j))$, which is at most
$O(d)+O(\nu-j)$ on an actual operand query.  If $\nu+v+h<0$, zero
is already a sufficiently accurate approximation after increasing
$h$, and no operand query is needed.  A fixed number of such
rescalings and clips preserves the bounds by telescoping.
Known non-power-of-two coefficients can instead be handled as
public normalized streams and bounded multiplication.
\end{proof}

\begin{lemma}[Finite comparisons with fixed margins]
\label{lem:arith-buffered-test}
Let $d\geq1$, let $s$ and $u$ have known magnitude bound $2^{\kappa d}$,
and suppose $u\geq\rho\geq2^{-\kappa d}$.  There is a deterministic
finite test that accepts every $s>2u$ and accepts only $s>u$.
It uses $O(d)$-place digit queries after normalization, with
intrinsic and suspended local lengths $O(d)$.  Its result is
independent of any subsequent numerical output digit position.
\end{lemma}

\begin{proof}
Compute finite rational approximations $\widetilde s,\widetilde u$
with individual errors at most $\rho/16$.  The known range and
lower bound on $\rho$ make $O(d)$ normalized digits sufficient.
Accept exactly when
$\widetilde s>\tfrac32\widetilde u$, using a fixed convention
at equality.  If $s>2u$, the left side minus the right side is
greater than $u/2-5\rho/32>0$.  If $s\leq u$, it is at most
$-u/2+5\rho/32<0$.  The finite comparison and all previous
arithmetic have the bounds already proved.  The test chooses its
precision from $d$ and $\rho$, not from an enclosing numerical
digit position.  With its structural starting digit position defined as zero, its
operand requests have $j=O(d)$; increasing $a$ in
Equation~\eqref{eq:arith-local-bound} pays also for the term $-bj$.
\end{proof}

\subsection{Proof of the numerical interface}
\label{subsec:arith-interface-proof}

\begin{proof}[Proof of Theorem~\ref{thm:arith-interface}]
The primitive constructions define deterministic integers $N_\nu$.
Lemma~\ref{lem:arith-residue-stream}
turns them into exact streams.  When composing primitives,
their operands are these exact streams; numerical approximation
errors are not accumulated along the recipe as if they were
statistical errors.  Instead, each primitive's deterministic
$N_\nu$ is an approximation to its own exact mathematical value.
The construction uses only residues of $N_\nu$ and $N_{\nu-1}$ and
a constant number of output code bits, which adds fixed
overshoot and fixed local overhead.

For the record bound, choose one common positive digit position
coefficient $b$ for the finite primitive and stage types.  In any
bound $L\leq a_0d+b_0(\nu-j)$ with $j-\nu\leq c_0d$, replacing
$b_0$ by $b\geq b_0$ is justified by replacing $a_0$ by
$a_0+(b-b_0)c_0$.  Thus the same $b$ can be used everywhere.
There are a bounded number of stages, each with overhead $O(d)$;
within a balanced sum or average the elementary overheads have total
$O(d)$ by Lemma~\ref{lem:arith-sums-scales}.  Along any path
the intermediate digit position differences telescope, giving
Equation~\eqref{eq:arith-export} and total overshoot $O(d)$.
The same sum, terminated at an internal arithmetic operation
rather than an operand call, gives the intrinsic bound explicitly
as follows.  If that operation has current digit position $u$, its own
intrinsic record has length at most $C(d+u)$, while the outer
saved portion has length at most $A d+b(\nu-u)$.  Choose the
common $b$ also to satisfy $b\geq C$.  Their sum is then at most
$(A+C)d+b\nu$, because $u\geq0$.  Increasing the stage constants
as above makes this same choice of $b$ valid in all local bounds.
The early-word case is already included in
Lemma~\ref{lem:arith-reciprocal}: those retained bits are part of
the local record and are paid for at the actual original-input
digit position by Equation~\eqref{eq:arith-early-word-payment}.

Structural validity branches are finite branches in this circuit
organization, with length at most its full intrinsic bound.
They ask the structural operand at digit position zero.  Comparisons with fixed margins
are separate recipes with structural starting digit position zero and
precision fixed by Lemma~\ref{lem:arith-buffered-test}; they are
not altered when a surrounding value is requested more accurately.
Thus no discontinuous comparison of exact reals at unbounded
precision is part of the interface.

All circuit constructions were given with padded domains and
index rules.  A saved gate choice selects an operand, a subtree,
a bit of a known-width word, or one of the finite primitive modes.
Scanning these choices from the root recovers the current
descriptor and digit places by integer arithmetic.  The
reciprocal's extra records additionally specify its early-prefix
formation state.  None of these computations refers to an
ancestor cursor address or requires an operand value to discover
the next gate.  Their integers have $O(B+d+\nu)$ bits; repeated
scans reuse one scratch area.  Lemma~\ref{lem:arith-gate-evaluation}
therefore applies at every Boolean portion of the recipe.

Finally, for a public coefficient $P/Q$, $Q>0$, producing a
place-$r$ approximation uses integer division of $2^rP$ by $Q$,
with signed floors when needed.  The inputs and intermediate
remainders have $O(B+r)$ bits.  Standard finite binary division
and multiplication use that much space, and extracting the
low residue or difference digit requires no larger word.
All operand bits in this computation are public, so it makes no
recursive graph call and can use shared scratch throughout.
For $r=O(B)$ its scratch bound is $O(B)$.  In contrast, any word
whose bits depend on recursive operands and must survive a new
operand request has already been charged to the local record.
\end{proof}

\begin{remark}[Use of digit positions]
\label{rem:arith-clocks}
Theorem~\ref{thm:arith-interface} is a finite-request statement
for every $\nu$.  Its exact-stream conclusion follows by considering
these finite requests separately, and does not require a machine
to choose structural domains or port orders at infinite precision.
For the actual top-level computation all reached digit positions will be
$O(B)$ by the allowance calculation in the next section.  Known
rate and field coefficients have $O(B)$-bit descriptions, and a
fixed rational generator probability raised to a walk length
$O(d)$ has numerator and denominator of $O(d)$ bits; these are
instances of the public-coefficient provision above.
\end{remark}
